\subsubsection{Hiện thực Instruction Register (IR)}
\begin{lstlisting}[language=Verilog, caption={Instruction Register}]
module instruction_reg (
    input             clk,      // Xung clock he thong
    input             rst,      // Reset dong bo 
    input             ld_ir,    // Tin hieu cho phep nap lenh tu Controller
    input      [31:0] data_in,  // Du lieu lenh tu Bus bo nho
    output reg [2:0]  opcode,   // 3-bit ma lenh 
    output reg [28:0] addr      // Dia chi bo nho di kem lenh
);

    // Instruction Register hoat dong dong bo theo canh len cua clock
    always @(posedge clk) begin
        if (rst) begin
            // Khoi tao gia tri ve 0 khi co tin hieu reset
            opcode <= 3'b000;
            addr   <= 29'b0;
        end
        else if (ld_ir) begin
            // Nap du lieu tu Bus vao thanh ghi khi ld_ir duoc kich hoat
            opcode <= data_in[31:29]; 
            addr   <= data_in[28:0];  
        end
    end

endmodule
\end{lstlisting}

\textbf{Mô tả logic xử lý:}
\begin{itemize}
    \item \textbf{Đồng bộ hóa tuyệt đối:} Khối \texttt{always} theo xung lên của clk. 
    Tín hiệu \texttt{rst} được kiểm tra ngay bên trong cấu trúc điều kiện của khối \texttt{always} thực hiện việc xóa dữ liệu (\texttt{opcode <= 0, addr <= 0}) theo xung clk (Reset đồng bộ).
    \item \textbf{Truyền dữ liệu:} Khi tín hiệu enable \texttt{ld\_ir} ở mức logic 0, 
    thanh ghi tự động duy trì trạng thái cũ. 
\end{itemize}